\section*{NeurIPS Paper Checklist}

\begin{enumerate}

\item {\bf Claims}
  \item[] Question: Do the main claims made in the abstract and introduction accurately reflect the paper's contributions and scope?
  \item[] Answer: \answerYes{}
  \item[] Justification: The abstract and introduction scope the paper to suppressive-workflow hazards under proxy-selected reporting, distinguish baseline-referenced workflow comparisons from within-trajectory selector/admissibility analyses, include the narrower CivilComments and ACSIncome support cases, and state the end-to-end Camelyon17 result without presenting it as selector-only causal evidence.

\item {\bf Limitations}
  \item[] Question: Does the paper discuss the limitations of the work performed by the authors?
  \item[] Answer: \answerYes{}
  \item[] Justification: The paper explicitly covers anchor dependence, the asymmetric cross-modal scope of the evidence (frozen-embedding text versus end-to-end vision), the heuristic nature of the guardrail, the stress-test status of the main P95 regime, and the fact that persistence is regime-dependent rather than universal.

\item {\bf Theory assumptions and proofs}
  \item[] Question: For each theoretical result, does the paper provide the full set of assumptions and a complete (and correct) proof?
  \item[] Answer: \answerYes{}
  \item[] Justification: The main text uses the local-lure condition, the baseline-referenced diagnostic-veto bound, and the transient-versus-persistent two-phase split only as compact organizing intuition, and Appendix~\ref{app:theory-bridge} gives the assumptions and proofs for each.

\item {\bf Experimental result reproducibility}
  \item[] Question: Does the paper fully disclose all the information needed to reproduce the main experimental results of the paper to the extent that it affects the main claims and/or conclusions of the paper (regardless of whether the code and data are provided or not)?
  \item[] Answer: \answerYes{}
  \item[] Justification: The paper specifies the datasets, training families, P95 distortion setting, selected-versus-fixed protocol, seed counts, held-out metrics, selector-dose checks, validation-accuracy selector sensitivity, leave-one-hospital-out selector diagnostics, the end-to-end GCE finetune selector audit, and diagnostic-veto sweep, and the supplement provides the exact commands, saved claim-bearing artifacts, and artifact-backed scripts used to regenerate the listed claim-bearing paper-facing figures and tables from those saved outputs. Full end-to-end reruns additionally require the external benchmark assets and the surrounding project environment.

\item {\bf Open access to data and code}
  \item[] Question: Does the paper provide open access to the data and code, with sufficient instructions to faithfully reproduce the main experimental results, as described in supplemental material?
  \item[] Answer: \answerYes{}
  \item[] Justification: The anonymized supplement includes the paper source, configs, exact commands, minimal dependencies, claim-bearing CSV artifacts, and the scripts that regenerate the listed paper-facing tables, figures, selector-dose diagnostics, validation-accuracy selector sensitivity, leave-one-hospital-out selector diagnostics, the end-to-end GCE finetune selector audit, and appendix seed-count / uncertainty summaries from saved outputs; benchmark datasets themselves are not redistributed and must be obtained under their original terms. This provides an open artifact-backed regeneration path for the reported claim-bearing results.

\item {\bf Experimental setting/details}
  \item[] Question: Does the paper specify all the training and test details (e.g., data splits, hyperparameters, how they were chosen, type of optimizer) necessary to understand the results?
  \item[] Answer: \answerYes{}
  \item[] Justification: The paper states the datasets, deployment splits, training families, P95 distortion setting, guardrail budgets, selected-versus-fixed protocol, seed counts, and evaluation metrics, and the reproducibility appendix lists the main optimizer, learning-rate, batch-size, epoch, and finetuning details used by the claim-bearing runs.

\item {\bf Experiment statistical significance}
  \item[] Question: Does the paper report error bars suitably and correctly defined or other appropriate information about the statistical significance of the experiments?
  \item[] Answer: \answerYes{}
  \item[] Justification: The main text surfaces paired bootstrap intervals for the principal loss / reliability / ECE harms on the three seed-matched hazard anchors, the appendix reports paired intervals for the Camelyon same-trajectory selector, selector-dose, validation-accuracy selector-sensitivity diagnostics, and the change from proxy-only to the diagnostic veto, and the main tables are built from explicit seed-matched summaries. The leave-one-hospital-out selector-only table is reported as a fold-level directional robustness check.

\item {\bf Experiments compute resources}
  \item[] Question: For each experiment, does the paper provide sufficient information on the computer resources (type of compute workers, memory, time of execution) needed to reproduce the experiments?
  \item[] Answer: \answerYes{}
  \item[] Justification: The reproducibility section states that the reported supervised experiments were run on a single NVIDIA GeForce RTX 5090 workstation, gives representative training hyperparameters, records the longest saved supporting run at roughly 12 GPU-hours total, and notes that the final paper assets are generated from saved repository artifacts by explicit scripts.

\item {\bf Code of ethics}
  \item[] Question: Does the research conducted in the paper conform, in every respect, with the NeurIPS Code of Ethics \url{https://neurips.cc/public/EthicsGuidelines}?
  \item[] Answer: \answerYes{}
  \item[] Justification: The work studies public benchmark datasets and evaluation suites, reports limitations and possible misuse, does not involve human-subject experimentation, and does not release high-risk models.

\item {\bf Broader impacts}
  \item[] Question: Does the paper discuss both potential positive societal impacts and negative societal impacts of the work performed?
  \item[] Answer: \answerYes{}
  \item[] Justification: The broader-impact section discusses the reliability value of the work, the risk of using objective shaping to improve surface metrics while hiding tail failures, and the importance of conservative reporting and deployment cross-checks in high-stakes settings.

\item {\bf Safeguards}
  \item[] Question: Does the paper describe safeguards that have been put in place for responsible release of data or models that have a high risk for misuse (e.g., pre-trained language models, image generators, or scraped datasets)?
  \item[] Answer: \answerNA{}
  \item[] Justification: The paper does not release a high-risk generative model, scraped dataset, or similarly misuse-prone asset.

\item {\bf Licenses for existing assets}
  \item[] Question: Are the creators or original owners of assets (e.g., code, data, models), used in the paper, properly credited and are the license and terms of use explicitly mentioned and properly respected?
  \item[] Answer: \answerYes{}
\item[] Justification: The supplement now includes \texttt{ASSET\_TERMS.md}, which credits the third-party datasets, benchmark suites, and code assets used in the supervised experiments, records the verified governing terms where available (including MIT, BSD-3-Clause, CC0, research-use terms, and public ACS microdata terms), and documents cases with ambiguous public notices. The submission redistributes only derived summaries and paper-facing artifacts, not raw images or benchmark shards.

\item {\bf New assets}
  \item[] Question: Are new assets introduced in the paper well documented and is the documentation provided alongside the assets?
  \item[] Answer: \answerYes{}
  \item[] Justification: The paper introduces paper-specific evaluation artifacts for the teacher-difficulty tail diagnostic on the vision anchors, and those artifacts are documented in Appendix~\ref{app:tail-bank-robustness} and the supplement through the bank-construction description, exact script names, paper-facing bankwise summaries, and artifact-backed regeneration commands. These are documented as reproducibility assets rather than as a standalone benchmark contribution.

\item {\bf Crowdsourcing and research with human subjects}
  \item[] Question: For crowdsourcing experiments and research with human subjects, does the paper include the full text of instructions given to participants and screenshots, if applicable, as well as details about compensation (if any)?
  \item[] Answer: \answerNA{}
  \item[] Justification: The reported work does not involve crowdsourcing or direct research with human participants.

\item {\bf Institutional review board (IRB) approvals or equivalent for research with human subjects}
  \item[] Question: Does the paper describe potential risks incurred by study participants, whether such risks were disclosed to the subjects, and whether Institutional Review Board (IRB) approvals (or an equivalent approval/review based on the requirements of your country or institution) were obtained?
  \item[] Answer: \answerNA{}
  \item[] Justification: The reported work does not involve human-subject experimentation.

\item {\bf Declaration of LLM usage}
  \item[] Question: Does the paper describe the usage of LLMs if it is an important, original, or non-standard component of the core methods in this research? Note that if the LLM is used only for writing, editing, or formatting purposes and does \emph{not} impact the core methodology, scientific rigor, or originality of the research, declaration is not required.
  \item[] Answer: \answerNA{}
  \item[] Justification: LLMs are not part of the method, experiment pipeline, or scientific contribution of this paper.

\end{enumerate}
